\subsection{Implementación de los cuantizadores (módulo \texttt{quantizers})}
\label{sec:implementation_quantizers}

\subsubsection{Clase \texttt{Quantizer}}

La interfaz de \texttt{TFMOT} \texttt{Quantizer} \cite{tensorflow_quantizer_api} define las operaciones matemáticas que transforman tensores durante la cuantización, la cual define dos métodos principales.

\begin{enumerate}
    \item \texttt{build}, que construye las variables de estado necesarias para el algoritmo de cuantización.
    \item \texttt{\_\_call\_\_}, que aplica la operación de cuantización a los tensores de entrada.
\end{enumerate}

El método \texttt{build} se llama una vez que se conoce la forma de los tensores de entrada, y permite inicializar las variables de estado. El método \texttt{\_\_call\_\_} se llama cada vez que se aplica la cuantización, y debe implementar la lógica de cuantización, esta es la que se coloca en el grafo computacional explicado en la sección \ref{sec:tensorflow_model_optimization}, para ejecutarse durante el \textit{forward pass}, las operaciones realizadas en este método deben tener definidas sus derivadas para permitir el cálculo de gradientes durante el \textit{backward pass}.

\subsubsection{Modificaciones a la clase \texttt{Quantizer}}
Para la implementación de cuantizadores con parámetros entrenables y gradientes personalizados, se realizaron algunas modificaciones sobre la clase \texttt{Quantizer} de \texttt{TFMOT}.

Para permitir definir gradientes personalizados, se define un método adicional \texttt{quantize} decorado con \texttt{@tf.custom\_gradient}, explicado en la sección \ref{sec:gradientes_personalizados}, donde se define la función de cuantización y su gradiente. El método \texttt{\_\_call\_\_} llama a \texttt{quantize} para aplicar la cuantización con el gradiente personalizado. Este mecanismo es el que se muestra en el fragmento de código \ref{code:quantizer_custom_gradient}.

\begin{lstlisting}[language=Python,label=code:quantizer_custom_gradient,caption=Definición de un cuantizador con gradiente personalizado.]
  class MyQuantizer(Quantizer):
    def __call__(self, inputs, training, weights, **kwargs):
      return self.quantize(inputs, **kwargs)

    @tf.custom_gradient
    def quantize(self, x, **kwargs):
      # Operacion de cuantizacion
      y = ... # Cuantizacion de x

      def grad(dy):
        # Gradiente personalizado
        return dy * ... # Derivada de la operacion de cuantizacion

      return y, grad
\end{lstlisting}

Para definir parámetros entrenables, se modificó el método \texttt{build} para declarar las variables de estado como tensores entrenables de Keras (\texttt{tf.Variable}). Esto permite que las variables de estado se entrenen junto con el modelo, y se actualicen durante el proceso de entrenamiento a partir de los gradientes calculados.

Por otro lado, los parámetros entrenables de los cuantizadores tiene restricciones específicas, por ejemplo los niveles y umbrales de la cuantización deben siempre estar ordenados, el parámetro que define el rango $\alpha$ debe ser siempre positivo, entre otros.

Para definir este comportamiento se implementaron restricciones en los parámetros entrenables utilizando la interfaz de \texttt{Keras} \texttt{Constraint} \cite{tensorflow_keras_constraint}, la cual define una operación de restricción que se aplica a los tensores entrenables durante el entrenamiento, como parte de la actualización de los mismos. Se implementaron las siguientes restricciones personalizadas en el módulo \texttt{constraints}:

\begin{itemize}
    \item \texttt{PositiveConstraint} que asegura que un tensor sea siempre positivo.
    \item \texttt{LevelConstraint} que asegura que los niveles sean siempre ordenados, esten dentro del rango de cuantización definido por $\alpha$ y que el primer nivel esta fijado al mínimo del rango $m$.
    \item \texttt{ThresholdConstraint} que asegura lo mismo que \texttt{LevelConstraint} pero fijando el ultimo nivel a el máximo del rango $M$.
\end{itemize}

%\TODO{Quizas agregar todo el codigo en apendice}
En el fragmento de código \ref{code:quantizer_constraints} se muestra la implementación de la restricción \texttt{PositiveConstraint}. Donde se utiliza la función \texttt{tf.clip\_by\_value} \cite{tensorflow_clip_by_value} para limitar los valores del tensor entre un valor mínimo (en este caso, un valor muy pequeño positivo definido por \texttt{tf.keras.backend.epsilon()} \cite{tensorflow_epsilon_api}) y el infinito positivo. Es necesario utilizar un valor mínimo mayor a cero para evitar divisiones por cero en los cálculos posteriores y asegurar la estabilidad numérica.

\begin{lstlisting}[language=Python,label=code:quantizer_constraints,caption=Definición de restricciones para parámetros entrenables.]
  class PositiveConstraint(tf.keras.constraints.Constraint):
    def __call__(self, w):
        return tf.clip_by_value(w, tf.keras.backend.epsilon(), np.inf)
\end{lstlisting}

Además los cuantizadores implementados exponen los inicializadores de los parámetros entrenables. Los parametros entrenables se agregan a la capa y se permite que el usuario pueda elegir como inicializar los parámetros entrenables de los cuantizadores a traves del argumento \texttt{initializer} \cite{tensorflow_initializers_api} del método \texttt{add\_weight} \cite{tensorflow_add_weight_api}. También se exponen los regularizadores en el módulo \texttt{regularizers}, para permitir que el usuario pueda elegir como regularizar los parámetros entrenables de los cuantizadores, también a traves del argumento \texttt{regularizer} \cite{tensorflow_regularizers_api} del método \texttt{add\_weight}.

La definición del método \texttt{build} se muestra en el fragmento de código \ref{code:quantizer_build}.

\begin{lstlisting}[language=Python,label=code:quantizer_build,caption=Definición del método build de un cuantizador con parámetros entrenables.]
  class MyQuantizer(Quantizer):
    def build(self, tensor_shape, name, layer):
      alpha = layer.add_weight(
          name.join("_alpha"),
          initializer=self.initializer,
          trainable=True,
          dtype=tf.float32,
          regularizer=self.regularizer,
          constraint=PositiveConstraint(),
        )
        return {"alpha": alpha}
\end{lstlisting}

\subsubsection{Cuantizador uniforme}

La implementación de la función de cuantización uniforme definida \eqref{eq:uniform_quantizer} se muestra en el fragmento de código \ref{code:uniform_quantizer}.

\begin{lstlisting}[language=Python,label=code:uniform_quantizer,caption=Definición de la función de cuantización uniforme.]
  def uniform_quantizer(x):
    clipped_x = tf.clip_by_value(x, levels[0], levels[-1])
    q = self.delta() * tf.math.floor(clipped_x / self.delta())
\end{lstlisting}

\TODO{Considerar agregar el histograma input/output de la función de cuantizacion}

Para implementar los gradientes se utilizan las definiciones de la sección \ref{sec:uniform_quantizer_gradient}, y se utilizan en la función \texttt{grad} dentro del método \texttt{quantize}.

El gradiente de la función de cuantización uniforme con respecto a la entrada $x$, $\frac{\partial q_{u}}{\partial x}$, definida \eqref{eq:uniform_quantizer_ste_signed} se implementa como se muestra en el fragmento de código \ref{code:uniform_quantizer_gradient_x}.

\begin{lstlisting}[language=Python,label=code:uniform_quantizer_gradient_x,caption=Definición del gradiente de la función de cuantización uniforme con respecto a la entrada $x$.]
  dq_dx = tf.where(
      tf.logical_and(
          tf.greater_equal(x, levels[0]),
          tf.less_equal(x, levels[-1]),
      ),
      upstream,
      tf.zeros_like(x),
  )
\end{lstlisting}

El gradiente de la función de cuantización uniforme con respecto al parámetro $\alpha$, $\frac{\partial q_{u}}{\partial \alpha}$, definida \eqref{eq:uniform_quantizer_dq_dalpha} se implementa como se muestra en el fragmento de código \ref{code:uniform_quantizer_gradient_alpha}. Donde \texttt{upstream} es el gradiente acumulado.

\begin{lstlisting}[language=Python,label=code:uniform_quantizer_gradient_alpha,caption=Definición del gradiente de la función de cuantización uniforme con respecto al parámetro $\alpha$.]
  dq_dalpha = tf.reduce_sum(q / alpha * upstream)
\end{lstlisting}

Es importante notar la utilización de \texttt{tf.reduce\_sum} en la implementación del gradiente con respecto a $\alpha$. Esto se debe a que $\alpha$ es un escalar, y la función de cuantización se aplica elemento a elemento sobre el tensor de entrada $x$. Por lo tanto, para calcular el gradiente total con respecto a $\alpha$, es necesario sumar los gradientes parciales de cada elemento del tensor de entrada.

\TODO{Mostrar entrenamiento, gif copado y evolucion comun}

\subsubsection{Cuantizador flexible}

La implementación de la función de cuantización flexible definida \eqref{eq:flexible_quantizer} se muestra en el fragmento de código \ref{code:flexible_quantizer}.

\begin{lstlisting}[language=Python,label=code:flexible_quantizer,caption=Definición de la función de cuantización flexible.]
  def flexible_quantizer(x):
    # Quantize the levels using uniform quantization
    qlevels = uniform_quantizer(self.levels)

    # Quantize input
    q = tf.zeros_like(x)
    for i in range(self.thresholds.shape[0] - 1):
        q = tf.where(
            tf.logical_and(
                tf.greater_equal(x, self.thresholds[i]),
                tf.less_equal(x, self.thresholds[i + 1]),
            ),
            qlevels[i],
            q,
        )
\end{lstlisting}

El gradiente de la función de cuantización flexible con respecto a la entrada $x$, $\frac{\partial q_{f}}{\partial x}$, es idéntico al del cuantizador uniforme, al ser un STE. Lo mismo ocurre con el gradiente con respecto a $\alpha$, $\frac{\partial q_{f}}{\partial \alpha}$.

El gradiente de la función de cuantización flexible con respecto a los niveles $l$, $\frac{\partial q_{f}}{\partial l}$, descripta \eqref{eq:flexible_quantizer_dq_dl} se implementa como se muestra en el fragmento de código \ref{code:flexible_quantizer_gradient_levels}.

\begin{lstlisting}[language=Python,label=code:flexible_quantizer_gradient_levels,caption=Definición del gradiente de la función de cuantización flexible con respecto a los niveles $l$.]
  bin_indices = tf.searchsorted(
      qlevels, tf.reshape(q, (-1,)), side="left"
  )
  q_one_hot = tf.one_hot(bin_indices, depth=qlevels.shape[0])
  dq_dlevel = tf.matmul(
                  tf.transpose(q_one_hot), tf.reshape(upstream, (-1, 1))
              )
              dq_dlevel = tf.reshape(dq_dlevel, (-1,))
\end{lstlisting}

En este caso, la dimensión de la variable de entrada $x$ no corresponde con la dimensión de la variable de los niveles $l$. Por ende, se utiliza la acumulación de los gradientes, para todos los valores de entrada $x_j$ que caen dentro del intervalo definido por los umbrales $t_i$ y $t_{i+1}$. Es decir, todas aquellas entradas que son cuantizadas al nivel $l_i$. Para esto, se utiliza \texttt{tf.searchsorted} para encontrar el índice del nivel correspondiente a cada valor cuantizado, luego se utiliza \texttt{tf.one\_hot} para crear una matriz one-hot con estos índices, y finalmente se utiliza \texttt{tf.matmul} para multiplicar la matriz one-hot transpuesta por el gradiente acumulado (\texttt{upstream}).

El gradiente de la función de cuantización flexible con respecto a los umbrales $t$, $\frac{\partial q_{f}}{\partial t}$, descripta \eqref{eq:flexible_quantizer_dq_dt} se implementa como se muestra en el fragmento de código \ref{code:flexible_quantizer_gradient_thresholds}.

\begin{lstlisting}[language=Python,label=code:flexible_quantizer_gradient_thresholds,caption=Definición del gradiente de la función de cuantización flexible con respecto a los umbrales $t$.]
  dq_dthresholds = tf.zeros_like(thresholds)
  for i in range(1, self.thresholds.shape[0] - 1):
      delta_y = qlevels[i - 1] - qlevels[i]
      delta_x = thresholds[i + 1] - thresholds[i - 1]

      # Only those associated with the 'x' values that
      # Fall within the range of the two borderline levels
      masked_upstream = tf.where(
          tf.logical_and(
              tf.greater_equal(x, self.thresholds[i - 1]),
              tf.less_equal(x, self.thresholds[i + 1]),
          ),
          upstream,
          tf.zeros_like(x),
      )

      update_value = (
          delta_y / delta_x * tf.reduce_sum(masked_upstream)
      )

      dq_dthresholds = tf.tensor_scatter_nd_update(
          dq_dthresholds, indices=[[i]], updates=[update_value]
      )
\end{lstlisting}

En este caso también se utiliza la acumulación de los gradientes, para todos los valores de entrada $x_j$ pues la dimensión de la variable de entrada $x$ no corresponde con la dimensión de la variable de los umbrales $t$. Para esto, se utiliza \texttt{tf.where} para crear una máscara que seleccione solo aquellos valores de $x$ que caen dentro del rango definido por los umbrales $t_{i-1}$ y $t_{i+1}$, luego se utiliza \texttt{tf.reduce\_sum} para sumar los gradientes correspondientes a estos valores, y finalmente se utiliza \texttt{tf.tensor\_scatter\_nd\_update} para actualizar el valor del gradiente correspondiente al umbral $t_i$.

\TODO{Mostrar entrenamiento, gif copado y evolucion comun}
